\chapter{Analysis: muons in the detector}
\label{ch:analysis}

\chaptergoals{10 of 11 in Part~II.}
{\cref{ch:production,ch:aegir,ch:digitisation}: the three files; the Python of
\cref{ch:team}.}
{find out which detectors the muons reach, with honest uncertainties, see
what the field does to them, and separate the two charges in the
spectrometer.}

There is no analysis framework to learn in this chapter. The files are
RNTuples of plain structs, so a Python script with PyROOT can read them
directly, and that is what we do.

\section{Reading the files}

The pattern from \cref{ex:openfile} is all you need: open a reader on the
RNTuple, create a typed \emph{view} per field, and call the view with an entry
number to get that entry's value as a C++ object. The scripts in
\file{code/analysis/} share a few small helpers in \file{shipfiles.py}: a
class that keeps readers and views together, the detector numbering, the
positions of the decay volume and the straw stations, and the normalisation
to protons on target.

\codefile[linerange={48-99},basicstyle=\ttfamily\footnotesize]{Python}{code/analysis/shipfiles.py}

\begin{pitfall}
Two habits of PyROOT with RNTuples are built into these helpers. A view
returns a reference into its own buffer, which the next call overwrites: copy
out what you need before reading the next entry. And tearing down readers
when Python exits can crash; the scripts therefore keep every reader alive and
leave through \texttt{os.\_exit} after flushing their output
(\texttt{sf.finish()}).
\end{pitfall}

\section{Which file for which question}

\begin{center}
\small
\begin{tabularx}{\linewidth}{X l}
\toprule
Question & File \\
\midrule
what leaves the hadron absorber, per proton & production (\cref{ch:production}) \\
per muon: which detectors did it reach & muon sample (\cref{ch:aegir}) \\
per event: which detectors saw anything & muon sample \\
detector response over all hits: resolutions, times & digitisation \\
\bottomrule
\end{tabularx}
\end{center}

The reason for the split is the shape of \pkg{Shannon}'s output
(\cref{sec:digilayout}): a stream of time-stamped hits, as the detector will
deliver, with no simulated-event number. Questions about one simulated event
or one muon are therefore answered from the simulation file, using simulated
hits as a stand-in for digitised ones. For counting which detectors were hit,
the position smearing does not matter; it would for reconstructed tracks,
timing or a realistic detector response.

\section{Fractions, not rates}

The muon sample of \cref{ch:aegir} is artificial: 100 muons of one charge,
flat in momentum, from one point at the back of the shield. What it can tell
you is what \emph{fraction} of such muons reaches each detector, not how many
muons per proton or per spill SHiP will see. Rates per proton come from the
plane of \cref{ch:production} or from a replay of a real production
(\cref{ex:ownreplay}); \file{muon\_background.py} normalises to protons only
when you give it \cmd{--pot}.

\begin{pitfall}
A fraction $f$ measured on $N$ muons has a binomial uncertainty
$\sqrt{f(1-f)/N}$: up to 5\,\% for $N=100$. Zero muons out of 100 means a
fraction below about 3\,\% (at 95\,\% confidence), not zero. Quote it that
way, and simulate more muons before reading anything into a small difference.
\end{pitfall}

\section{Where the muons end up}

\file{code/analysis/muon\_background.py} looks at every primary muon (a
\texttt{SimParticle} with \field{parent\_id}~$=0$ and $|\field{pdg\_code}|=13$)
and finds the detectors in which \emph{the muon itself} left hits, and the
straw stations it crossed. Separately, it counts the events with any hit in
each detector, whoever made it: a muon that only passes close to a detector
can still send secondaries into it.

\begin{exercise}{Run the analysis}
\label{ex:muonbackground}
In the \pkg{aegir} environment, in \file{work/}:
\begin{shell}
python $CODE/analysis/muon_background.py muon_gun.root
\end{shell}
It prints two tables and writes three PDFs (\file{muon\_background\_*.pdf}).
\begin{enumerate}
  \item What fraction of the muons reaches the UBT, the straws, all four
        straw stations, the timing detector and the calorimeter? Give each
        with its uncertainty.
  \item Which muons do not cross all four stations?
  \item Compare the two tables: why can an event have hits in a detector
        that its muon did not reach?
  \item \file{\_p.pdf}: does the fraction reaching the straws depend on the
        momentum?
  \item \file{\_station1\_xy.pdf}: how large is the spot of the muons in
        the first straw station, and does that agree with 20\,mrad over the
        distance from $z=27$\,m?
\end{enumerate}
\end{exercise}

\begin{checkpoint}
Two tables whose first line counts 100 primary muons, with a
\texttt{fraction} column, and three PDFs.
\end{checkpoint}

\begin{solution}
\begin{enumerate}
  \item Nearly all: at 20\,mrad or less the muons stay inside the decay
        volume. In the dry run, 100\,\% reached the UBT, 97\,\% the straws,
        95\,\% all four stations, the timing detector and the calorimeter,
        each $\pm2\,\%$, and 5\,\% the SBT ($\pm2\,\%$).
  \item The ones at the largest angles, which leave the acceptance before the
        last stations or cross the vessel wall into the SBT. Compare
        the 5\,\% that miss a station with the 5\,\% in the SBT: check in the
        script whether they are the same muons.
  \item Secondaries: a $\delta$-ray, a photon from bremsstrahlung or a hadron
        from a muon-nucleus interaction reaches a detector the muon itself
        misses. Such hits are why a veto has to be studied with the full
        detector simulation, not with the muon trajectories alone.
  \item Hardly, above a few GeV: the muons are too hard to be stopped, and
        the field of the magnet moves the soft ones by centimetres to
        decimetres only.
  \item At most $0.02\times(84.07-27)\,\mathrm{m}\approx1.1$\,m from the
        axis, and uniform in angle, so denser towards the centre.
\end{enumerate}
\end{solution}

\section{What the field does}

The spectrometer magnet between straw stations 2 and 3 bends the muons in the
horizontal plane; the muon shield, behind which the gun starts, no longer acts
on them. The sample without field (\cref{ex:nofield}) shows the difference.

\begin{exercise}{Field on, field off}
If you made \file{muon\_gun\_nofield.root}, run the analysis on it too:
\begin{shell}
python $CODE/analysis/muon_background.py muon_gun_nofield.root \
    --output nofield
\end{shell}
Compare the two muon tables. Do the fractions change, within their
uncertainties? Compare \file{muon\_background\_station1\_xy.pdf} with
\file{nofield\_station1\_xy.pdf}: why do they look alike?
\end{exercise}

\begin{solution}
The fractions agree within their uncertainties: the field changes where a
muon crosses stations 3 and 4 by centimetres to decimetres, small compared
with the size of the detectors. Station~1 is upstream of the magnet, so its
hit maps differ only by Geant4's randomness. The field does not decide
\emph{whether} a muon is seen; it decides where, which is what measures its
momentum and charge. \tbd{the two tables of the dry run}
\end{solution}

\begin{exercise}[colback=white]{Optional: the two charges}
With the $\mu^+$ sample of \cref{ex:mup}: for every muon that crosses all
four straw stations, take the mean $x$ and $z$ of its hits in each station
and compute the change of its slope in $x$ across the magnet,
$\Delta\theta_x = \frac{x_4-x_3}{z_4-z_3} - \frac{x_2-x_1}{z_2-z_1}$.
Histogram it separately for $\mu^+$ and $\mu^-$. Do the two charges go to
opposite sides? Repeat on the sample without field.
\end{exercise}

\begin{solution}
A solution is \file{code/analysis/solutions/spectrometer\_charge.py}:
\begin{shell}
python $CODE/analysis/solutions/spectrometer_charge.py \
    muon_gun.root muon_gun_mup.root
python $CODE/analysis/solutions/spectrometer_charge.py \
    muon_gun_nofield.root --output nofield
\end{shell}
With field the two distributions lie on opposite sides of zero, and the kink
falls with momentum as $\Delta\theta_x \approx 0.3\int B\,\mathrm{d}l / p$
(\cref{ch:services}), so the soft muons give the largest kinks. Without field
the distribution is centred on zero, with the width that multiple scattering
gives.
\end{solution}

\section{Going further}

\begin{exercise}[colback=white]{Optional: from a hit stream back to events}
Using the digitised file, histogram the times of all SBT hits. Could you
separate the simulated events by time? What would the digitisation have to
record so that a Monte Carlo study can connect digitised hits to the
simulated event they came from? Bring your answer to the closing discussion.
\end{exercise}

\begin{exercise}[colback=white]{Optional: more muons}
Simulate 1000 muons with another seed and another \texttt{output} name
(\cmd{--tla-code events=1000 --tla-code seed=2}) and run the analysis on
them. How do the fractions and their uncertainties change? How many protons
on target would it take to get 1000 muons behind the shield from the real
production of \cref{ch:production}?
\end{exercise}

\begin{solution}
The fractions agree with the 100-muon sample within its uncertainties, and
the uncertainties shrink by $\sqrt{10}$. For the protons: the muons per proton
at the plane from \cref{ex:planesummary}, times the fraction that gets
through the shield (\cref{ex:ownreplay}), gives the muons behind the shield
per proton; 1000 divided by it is the number of protons, which is why SHiP's
background studies use dedicated large productions, with event weights and
biasing.
\end{solution}
